\documentclass[aspectratio=169]{beamer}

\usetheme{Madrid}
\usepackage[utf8]{inputenc}
\usepackage[T1]{fontenc}
\usepackage{amsmath, amssymb}
\usepackage{booktabs}
\usepackage{array}
\usepackage{xcolor}
\usepackage{microtype}

\setbeamertemplate{navigation symbols}{}
\setbeamertemplate{footline}[frame number]

\definecolor{mainblue}{RGB}{0, 51, 102}
\setbeamercolor{structure}{fg=mainblue}

\title[Putujući lopov]{Putujući lopov}
\subtitle{Metaheuristički pristupi}
\author{Anđela Spasić \and Matija Radulović}
\institute{Matematički fakultet, Univerzitet u Beogradu}
\date{Jun 2026}

\begin{document}

\begin{frame}
  \titlepage
\end{frame}

\begin{frame}{Sadržaj}
  \tableofcontents
\end{frame}

% =========================================================
\section{Problem}
% =========================================================

\begin{frame}{Putujući lopov (TTP)}
  \begin{block}{Definicija}
    Lopov obilazi $n$ gradova tačno jednom i može da uzme predmete koji se
    nalaze u gradovima. Plaća rentu $R$ po jedinici vremena.
    Cilj: maksimizovati profit minus rentu.
  \end{block}

  \bigskip
  \begin{itemize}
    \item Kombinacija \textbf{TSP} i \textbf{0/1 ranca} --- oba NP-teška
    \item Svaki uzeti predmet \emph{usporava} lopova do kraja ture
    \item \textbf{NP-težak} (Bonyadi, Michalewicz, Barone --- IEEE CEC 2013)
  \end{itemize}

  \bigskip
  \begin{alertblock}{Zašto nije trivijalno spajanje}
    Optimalna tura za TSP $\neq$ optimalna tura za TTP ---
    problemi se moraju rešavati zajedno.
  \end{alertblock}
\end{frame}

% ---------------------------------------------------------
\begin{frame}{Matematička formulacija}

  \textbf{Brzina} pri tekućoj težini $w$:
  \[
    v(w) = v_{\max} - w \cdot \frac{v_{\max} - v_{\min}}{W}
  \]

  \textbf{Vreme prelaska} između $c_k$ i $c_{k+1}$:
  \[
    t_k = \frac{d(c_k,\, c_{k+1})}{v(w_k)}
  \]

  \textbf{Ciljna funkcija} (maksimizovati):
  \[
    f = \sum_{i \in \text{odabrani}} p_i \;-\; R \cdot \sum_{k=1}^{n} t_k
  \]
\end{frame}

% ---------------------------------------------------------
\begin{frame}{Instance}
  \begin{center}
  \begin{tabular}{lrrll}
    \toprule
    \textbf{Instanca} & \textbf{$n$} & \textbf{$m$} & \textbf{Tip predmeta} & \textbf{Namena} \\
    \midrule
    tiny5  &  5 &  4 & ---                    & brute force provera \\
    tiny7  &  7 &  6 & ---                    & brute force provera \\
    tiny9  &  9 &  8 & ---                    & brute force provera \\
    eil51\_uncorr & 51 & 50 & nekorelisani    & glavna evaluacija \\
    eil51\_corr   & 51 & 50 & jako korelisani & glavna evaluacija \\
    pr76\_uncorr  & 76 & 75 & nekorelisani    & veća instanca \\
    \bottomrule
  \end{tabular}
  \end{center}

  \medskip
  {\small Polyakovskiy et al., GECCO 2014}
\end{frame}

% =========================================================
\section{Referentna rešenja}
% =========================================================

\begin{frame}{Brute Force}
  \begin{itemize}
    \item Prolazi sve $(n-1)! \times 2^m$ kombinacije ture i pakovanja
    \item Garantovano optimalno rešenje
    \item \textbf{Složenost:} $\mathcal{O}\!\left((n-1)!\cdot 2^m\right)$
  \end{itemize}

  \bigskip
  \begin{block}{Izvodljivost}
    \begin{tabular}{lrl}
      tiny7 ($n=7, m=6$)    & $720 \times 64 = 46\,080$    & \textcolor{green!60!black}{$\checkmark$ izvodljivo} \\
      eil51 ($n=51, m=50$)  & $50!\times 2^{50}\approx 10^{80}$ & \textcolor{red}{$\times$ nemoguće}
    \end{tabular}
  \end{block}
\end{frame}

% ---------------------------------------------------------
\begin{frame}{S5 --- Pohlepna deterministička heuristika}
  \begin{enumerate}
    \item Nearest-neighbor tura od grada 0 \hfill $\mathcal{O}(n^2)$
    \item Pohlepno pakovanje po skoru:
      \[
        \text{score}(i) = \frac{p_i}{w_i \cdot \text{preostalo\_rastojanje}(c_i)}
      \]
      \hfill $\mathcal{O}(m \log m)$
    \item Iterativno poboljšanje pakovanja (bit-flip) \hfill $\mathcal{O}(m^2)$
  \end{enumerate}

  \bigskip
  \textbf{Ukupna složenost:} $\mathcal{O}(n^2 + m^2)$ \quad --- \quad
  Donja referentna tačka za sve algoritme.
\end{frame}

% =========================================================
\section{Simulirano kaljenje}
% =========================================================

\begin{frame}{SA --- Simulirano kaljenje}
  \textbf{Parametri:} 50\,000 iteracija, $T_0 = 500$, $c = 0{,}9995$, $T_{\min} = 1$

  \bigskip
  \textbf{Potezi po iteraciji:}
  \begin{itemize}
    \item 40\% --- 2-opt preokret segmenta ture
    \item 30\% --- premeštanje jednog grada na drugu poziciju
    \item 40\% --- bit-flip jednog predmeta u rancu
  \end{itemize}

  \bigskip
  \textbf{Bolmanovo prihvatanje lošijeg rešenja:}
  \[
    P(\Delta) = e^{\,\Delta / T}, \qquad T \leftarrow c \cdot T
  \]
  \medskip
  \textbf{Složenost:} $\mathcal{O}(T_{\max} \cdot (n + m))$
\end{frame}

% ---------------------------------------------------------
\begin{frame}{SA Improved}
  \textbf{Inicijalizacija:} pohlepna tura $\to$ 2-opt $\to$ or-opt $\to$ iterativno pakovanje

  \bigskip
  \textbf{Auto-kalibracija temperature} (ciljano 80\% prihvatanje na startu):
  \[
    T_0 = \frac{-\,\overline{|\Delta|}}{\ln 0{,}8}
  \]

  \bigskip
  \textbf{Po iteraciji:}
  \begin{itemize}
    \item $n/2$ poteza za turu --- pohlepno prihvatanje (samo poboljšanje)
    \item $m/2$ flip-ova za ranac --- SA prihvatanje
  \end{itemize}

  \medskip
  \textbf{Složenost:} $\mathcal{O}(T_{\max} \cdot (n + m))$
\end{frame}

% =========================================================
\section{Mravlja kolonija}
% =========================================================

\begin{frame}{ACO --- Konstrukcija ture}
  \begin{tabular}{lll}
    $n_{\text{ant}} = n$    & broj mrava         & jednak broju gradova \\
    $\alpha = 1$            & uticaj feromona    & koliko mravi prate već poznate puteve \\
    $\beta = 2$             & uticaj heuristike  & privlačnost bliskih gradova ($1/d$) \\
    $\rho = 0{,}02$         & isparavanje        & brzina zaboravljanja starih puteva \\
    $p_{\text{best}} = 0{,}05$ & ciljana $p$     & koristi se za izračun $\tau_{\min}$ \\
  \end{tabular}

  \medskip
  \textbf{Verovatnoća izbora sledećeg grada:}
  \[
    P(i \to j) = \frac{\tau(i,j)^\alpha \cdot \eta(i,j)^\beta}
                       {\displaystyle\sum_{l \notin \text{posećeni}} \tau(i,l)^\alpha \cdot \eta(i,l)^\beta},
    \qquad \eta(i,j) = \frac{1}{d(i,j)}
  \]

  \bigskip
  Svaka mravi gradi kompletnu turu, zatim se pohlepno pakuje ranac.
\end{frame}

% ---------------------------------------------------------
\begin{frame}{ACO --- MMAS mehanizmi}
  \textbf{Granice feromona} (Stützle \& Hoos, 2000):
  \[
    \tau_{\max} = \frac{1}{\rho \cdot L_{\text{best}}}
    \quad \text{\small(dodati = uklonjeni feromoni --- stabilno stanje)}
  \]
  \[
    p = p_{\text{best}}^{1/n}
    \quad \text{\small($p^n = p_{\text{best}}$: verovatnoća biranja iste ivice kroz celu turu)}
  \]
  \[
    \tau_{\min} = \tau_{\max} \cdot \frac{1-p}{(\bar{k}-1)\cdot p},
    \quad \bar{k} = \tfrac{n}{2}
    \quad \text{\small(procena neposećenih gradova po koraku)}
  \]
  {\small $\tau_{\min}$: vrednost pri kojoj $\tau_{\max}$ ivica pobedi naspram $(\bar{k}-1)$ lošijih}

  \medskip
  \textbf{Evaporacija i depozit:}
  \[
    \tau(a,b) \;\leftarrow\; (1-\rho)\,\tau(a,b) \;+\; \frac{1}{L_{\text{deposit}}}
  \]
  {\small $\tau_{\max}$ sprečava dominaciju jednog puta; $\tau_{\min}$ sprečava zanemarivanje puteva}

  \medskip
  \begin{itemize}
    \item \textbf{Stagnacija:} reset $\tau \leftarrow \tau_{\max}$ posle 100 iteracija bez poboljšanja
    \item \textbf{Depozit:} iter-best (prva polovina) $\to$ global-best (druga polovina)
  \end{itemize}

  \textbf{Složenost:} $\mathcal{O}(\text{iter} \cdot (n^3 + m^2))$
\end{frame}

% =========================================================
\section{Genetski algoritam}
% =========================================================

\begin{frame}{GA --- Osnovna verzija}
  \textbf{Parametri:} populacija 50, $cr = 0{,}8$, $mr = 0{,}1$, turnir $k=3$, elitizam, 500 generacija

  \bigskip
  \begin{itemize}
    \item \textbf{Ukrštanje ture:} OX (order crossover) --- čuva relativni redosled gradova
    \item \textbf{Ukrštanje ranca:} uniformno, popravka kapaciteta
    \item \textbf{Mutacija:} 2-opt preokret / bit-flip predmeta
  \end{itemize}

  \bigskip
  \textbf{Složenost:} $\mathcal{O}(\text{iter} \cdot \text{pop} \cdot (n + m))$
\end{frame}

% ---------------------------------------------------------
\begin{frame}{GA Improved}
  \textbf{Parametri:} populacija 30, $cr = 0{,}85$, $mr = 0{,}15$, steady-state, 1000 iteracija

  \bigskip
  \begin{itemize}
    \item \textbf{Ukrštanje:} EAX (edge assembly crossover) ---
      čuva dobre grane iz oba roditelja kroz AB-cikluse
    \item \textbf{Po potomku:} 2-opt + or-opt + iterativno pakovanje
    \item \textbf{Inicijalizacija:} svi jedinke pohlepne ture sa različitim startnim gradovima
  \end{itemize}

  \bigskip
  \textbf{Složenost:} $\mathcal{O}(\text{iter} \cdot (n^2 + m^2))$
\end{frame}

% =========================================================
\section{Optimizator sivog vuka}
% =========================================================

\begin{frame}{GWO --- Pomeranje vukova}
  \textbf{Hijerarhija:} $\alpha$ (best), $\beta$, $\delta$ vode --- omega vukovi se pomeraju ka svakom vodiču zasebno, svaki sa nezavisnim izvlačenjem $r_1$.

  \medskip
  \textbf{Parametar istraživanja} (opada $2 \to 0$):
  \[
    a(t) = 2 - 2\,\frac{t}{T_{\max}}, \qquad
    |A| = |2\,a\,r_1 - a|, \quad r_1 \sim U(0,1)
  \]

  \textbf{Pomeranje ture} (ponavlja se ka $\alpha$, $\beta$, $\delta$ redom):
  \begin{itemize}
    \item $|A| \geq 1$: visoko $a$, rano u pretrazi --- nasumični 2-opt preokret (istraživanje, ignorisati vodiča)
    \item $|A| < 1$: nisko $a$, kasno u pretrazi --- primeni swap sekvence ka vodiču (iskorišćavanje):
      \[
        n_{\text{swap}} = \lfloor (1 - |A|) \cdot |\text{swaps}| \rfloor
      \]
  \end{itemize}
\end{frame}

% ---------------------------------------------------------
\begin{frame}{GWO --- Ažuriranje ranca}
  \textbf{Kopiranje bita od svakog vodiča} (nezavisno za $\alpha$, $\beta$, $\delta$):
  \[
    \text{copy\_prob} = \max(0,\; 1 - |A|)
  \]
  {\small $|A| \to 0$: kopiraj sve od vodiča \quad $|A| \to 1$: ne kopiraj ništa}

  \bigskip
  \textbf{Glasanje većine trojice vodiča} (bit je 1 ako $\geq 2$ od 3 vodiča kažu 1):
  \[
    x_k = \mathbf{1}\!\left[x_k^{(\alpha)} + x_k^{(\beta)} + x_k^{(\delta)} \geq 2\right]
  \]

  \begin{itemize}
    \item Glasanje može da prekorači kapacitet --- popravka izbacivanjem najlošijih predmeta
    \item \textbf{Inicijalizacija:} 90\% nasumično, 10\% perturbovano pohlepno rešenje
  \end{itemize}

  \medskip
  \textbf{Parametri:} $n_{\text{wolves}} = n$,\; 500 iteracija \quad
  \textbf{Složenost:} $\mathcal{O}(\text{iter} \cdot n \cdot (n + m))$
\end{frame}

% =========================================================
\section{Rezultati}
% =========================================================

\begin{frame}{Pregled složenosti}
  \begin{center}
  \renewcommand{\arraystretch}{1.4}
  \begin{tabular}{lll}
    \toprule
    \textbf{Algoritam} & \textbf{Složenost po iteraciji} & \textbf{Ukupna složenost} \\
    \midrule
    Brute Force  & ---                            & $\mathcal{O}((n-1)!\cdot 2^m)$ \\
    S5           & --- (bez iteracija)            & $\mathcal{O}(n^2 + m^2)$ \\
    SA           & $\mathcal{O}(n + m)$           & $\mathcal{O}(T_{\max}(n+m))$ \\
    ACO (MMAS)   & $\mathcal{O}(n^3 + m^2)$      & $\mathcal{O}(\text{iter}(n^3+m^2))$ \\
    GA           & $\mathcal{O}(\text{pop}(n+m))$ & $\mathcal{O}(\text{iter}\cdot\text{pop}(n+m))$ \\
    GA Improved  & $\mathcal{O}(n^2 + m^2)$      & $\mathcal{O}(\text{iter}(n^2+m^2))$ \\
    GWO          & $\mathcal{O}(n(n+m))$         & $\mathcal{O}(\text{iter}\cdot n(n+m))$ \\
    \bottomrule
  \end{tabular}
  \end{center}
\end{frame}

% ---------------------------------------------------------
\begin{frame}{Rezultati --- tiny7 ($n=7, m=6$)}
  \begin{center}
  \renewcommand{\arraystretch}{1.3}
  \begin{tabular}{lrrrr}
    \toprule
    \textbf{Algoritam} & \textbf{Najbolji} & \textbf{Prosek} & \textbf{Najgori} & \textbf{Vreme (s)} \\
    \midrule
    Brute Force          & \textbf{115.23} & ---    & ---    & 0.25 \\
    \midrule
    S5                   &  65.78 &  65.78 &  65.78 & 0.00 \\
    GWO                  & 110.40 &  90.44 &  50.54 & 0.03 \\
    GA Improved          & 101.70 & 101.70 & 101.70 & 0.09 \\
    GA                   & 110.40 & 109.91 & 108.95 & 0.30 \\
    SA                   & 115.23 & 111.53 & 108.95 & 0.11 \\
    ACO (MMAS)           & 115.23 & 115.23 & 115.23 & 0.09 \\
    \textbf{SA Improved} & \textbf{115.23} & \textbf{115.23} & \textbf{115.23} & 0.42 \\
    \bottomrule
  \end{tabular}
  \end{center}

  \smallskip
  {\footnotesize 3 pokretanja. Optimum (Brute Force) = 115.23.}
\end{frame}

% ---------------------------------------------------------
\begin{frame}{Konvergencija}
  \begin{center}
    \textit{[Grafik konvergencije --- eil51\_corr, svi algoritmi]}
  \end{center}

  \begin{itemize}
    \item SA --- glatka konvergencija kroz sitne poteze
    \item ACO --- skokovi vidljivi pri resetovima na stagnaciju
    \item GA Improved --- brže od GA zahvaljujući EAX i lokalnoj pretrazi
    \item GWO --- rano stagnira
  \end{itemize}
\end{frame}

% =========================================================
\section{Zaključak}
% =========================================================

\begin{frame}{Zaključak}
  \begin{block}{Rangiranje (eil51 i pr76)}
    \textbf{SA Improved} $>$ SA $>$ ACO $>$ GA Improved $>$ GA $>$ GWO $>$ S5
  \end{block}

  \bigskip
  \begin{itemize}
    \item \textbf{SA:} pretraga jednog rešenja sa sitnim potezima prirodno odgovara TTP-u
    \item \textbf{ACO:} lokalna pretraga + explore$\to$exploit depozit značajno pomažu
    \item \textbf{GA Improved:} EAX i lokalna pretraga na potomku nadoknađuju manju populaciju
    \item \textbf{GWO:} swap sekvence ne uspevaju da uhvate finu strukturu TTP-a
  \end{itemize}

  \bigskip
  \textbf{Moguća poboljšanja:} LKH za ture, dinamičko programiranje za ranac pri fiksnoj turi.
\end{frame}

% ---------------------------------------------------------
\begin{frame}{Literatura}
  \footnotesize
  \begin{itemize}
    \item Bonyadi, M.R., Michalewicz, Z., Barone, L. (2013).
      \textit{The Travelling Thief Problem: The First Step in the Transition
      from Theoretical Problems to Realistic Problems.} IEEE CEC.

    \item Polyakovskiy, S. et al. (2014).
      \textit{A Comprehensive Benchmark Set and Heuristics for the
      Traveling Thief Problem.} GECCO.

    \item Stützle, T., Hoos, H.H. (2000).
      \textit{MAX--MIN Ant System.}
      Future Generation Computer Systems, 16(8), 889--914.

    \item Mirjalili, S., Mirjalili, S.M., Lewis, A. (2014).
      \textit{Grey Wolf Optimizer.}
      Advances in Engineering Software, 69, 46--61.

    \item Kirkpatrick, S., Gelatt, C.D., Vecchi, M.P. (1983).
      \textit{Optimization by Simulated Annealing.}
      Science, 220(4598), 671--680.

    \item Nagata, Y., Kobayashi, S. (2013).
      \textit{A Powerful Genetic Algorithm Using Edge Assembly Crossover
      for the Traveling Salesman Problem.} INFORMS Journal on Computing.
  \end{itemize}
\end{frame}

\end{document}
